% 329_nthresh_Skalenanalyse_De_ch.tex
% Dok. 329, FFGFT-Korpus, August 2026 — Kapitel-Datei
% Skalenanalyse der Stabilitätsschwelle n_thresh

% --- Lokale Makros ---
\providecommand{\K}{}\renewcommand{\K}{\textbf{[K]}}
\providecommand{\B}{}\renewcommand{\B}{\textbf{[B]}}
\renewcommand{\S}{\textbf{[S]}}
\providecommand{\X}{}\renewcommand{\X}{\textbf{[X]}}
\providecommand{\lP}{}\renewcommand{\lP}{\ell_{\mathrm{P}}}
\providecommand{\mP}{}\renewcommand{\mP}{m_{\mathrm{P}}}
\providecommand{\TP}{}\renewcommand{\TP}{T_{\mathrm{P}}}
\providecommand{\Ebit}{}\renewcommand{\Ebit}{E_{\mathrm{bit}}}
\providecommand{\mbit}{}\renewcommand{\mbit}{m_{\mathrm{bit}}}
\providecommand{\nthr}{}\renewcommand{\nthr}{n_{\mathrm{thresh}}}
\providecommand{\Mcoll}{}\renewcommand{\Mcoll}{M_{\mathrm{coll}}}

\definecolor{ffgftblue}{RGB}{30,90,160}
\definecolor{proofgreen}{RGB}{0,110,50}
\definecolor{warnviolet}{RGB}{100,0,120}
\tcbset{
  base/.style={breakable,enhanced,boxrule=0.6pt,arc=3pt,
    fonttitle=\bfseries\small,coltitle=white,top=4pt,bottom=4pt,left=6pt,right=6pt},
  ffgftbox/.style={base,colback=blue!5,colframe=ffgftblue,colbacktitle=ffgftblue},
  proofbox/.style={base,colback=green!5,colframe=proofgreen,colbacktitle=proofgreen},
  warnbox/.style={base,colback=violet!6,colframe=warnviolet,colbacktitle=warnviolet}
}

% ============================================================
\section*{Zusammenfassung}

Douglas Matzkes Hyperbit-Framework (IPI 2026) leitet eine
Stabilitätsschwelle
\begin{equation}
  \nthr = \frac{2\pi}{\sqrt{2}\,\ln 2} = 6{,}4097~\text{Bit}
\end{equation}
her und deutet sie als universelle Zahl: $\mathrm{Cl}(6)$ trägt
$6{,}000$~Bit und liegt damit $0{,}41$~Bit unterhalb der Kollapsschwelle --
Materie ist stabil, weil die kleinste Fermion-Algebra gerade eben nicht
kollabiert. Dok.~326 hatte den Verdacht formuliert, dass diese Zahl an
der Setzung einer bestimmten Temperatur hängt. Dieses Dokument prüft das
numerisch und beantwortet es abschließend.

\begin{tcolorbox}[ffgftbox, title={Hauptresultat \B}]
  Mit der systemabhängigen Bit-Energie der FFGFT
  $\Ebit(L) = \hbar c/L$ (Dok.~257) folgt exakt
  \begin{equation}
    \nthr(L) = \frac{L}{\sqrt{2}\,\lP}
    \label{eq:main}
  \end{equation}
  -- \textbf{linear in der Skala}, kein invarianter Zahlenwert.
  Matzkes $6{,}41$ entspricht exakt der Skala
  \begin{equation}
    L_{\mathrm{Matzke}} = \frac{2\pi\,\lP}{\ln 2} = 9{,}0647\,\lP .
  \end{equation}
  Der \enquote{universelle} Bitwert \emph{ist} die FFGFT-Bit-Energie,
  ausgewertet bei rund neun Plancklängen.
  Die Universalität ist eine verdeckte Skalenwahl.
\end{tcolorbox}

Was dabei \emph{unangetastet} bleibt, ist der belastbare Kern von Matzkes
Argument: Die Kollapsbedingung Compton~$=$~Schwarzschild ergibt
$\Mcoll = \mP/\sqrt{2}$ rein geometrisch -- ohne jeden Bezug auf Bits,
Landauer oder eine Temperatur. Dieser Teil steht unabhängig \B.

Prüfskript: \texttt{326\_nthresh\_skalenabhaengigkeit.py} --
12 Assertionen, alle bestanden.

% ============================================================
\section{Die Frage}

Dok.~326 arbeitet die zentrale Divergenz zwischen Matzkes Framework und
FFGFT heraus: universelle vs.\ systemabhängige Bitwerte.
Der dort formulierte Verdacht lautete, dass Matzkes Stabilitätsschwelle
\enquote{an der Setzung $T = \TP$ hängt} und bei systemabhängiger
Bit-Energie selbst systemabhängig wäre -- eingestuft als \S.

Diese Frage ist numerisch entscheidbar, und zwar vollständig:
Man rechnet Matzkes Kette nach, extrahiert die skalenunabhängigen
Anteile, ersetzt den Bitwert durch die FFGFT-Bit-Energie und sieht,
ob eine Konstante übrig bleibt.

% ============================================================
\section{Matzkes Kette}

\subsection{Die Landauer-Brücke}

Die angesetzte Temperatur ist nicht die Plancktemperatur selbst,
sondern die \textbf{KMS-Temperatur} der Planckskala:
\begin{equation}
  T_{\mathrm{KMS}} = \frac{\TP}{2\pi} = 2{,}2549\times10^{31}~\mathrm{K}.
\end{equation}
Das $2\pi$ stammt aus der KMS-Periodizität -- derselben, die in
Dok.~325 die Hawking- und Unruh-Temperatur erzeugt.
Landauer bei dieser Temperatur:
\begin{equation}
  \Ebit = k_B T_{\mathrm{KMS}} \ln 2
  \quad\Longrightarrow\quad
  \mbit = \frac{\mP \ln 2}{2\pi} = 0{,}1103\,\mP .
\end{equation}
Numerisch reproduziert (Skript~[1], Verhältnis $1{,}000000000$). \B

\subsection{Die Kollapsbedingung}

Compton-Wellenlänge gleich Schwarzschild-Radius:
\begin{equation}
  \frac{\hbar}{Mc} = \frac{2GM}{c^2}
  \quad\Longrightarrow\quad
  \Mcoll = \sqrt{\frac{\hbar c}{2G}} = \frac{\mP}{\sqrt{2}}
         = 1{,}5390\times10^{-8}~\mathrm{kg}.
\end{equation}

\begin{tcolorbox}[proofbox, title={Der skalenunabhängige Kern \B}]
  $\Mcoll = \mP/\sqrt{2}$ enthält \textbf{keinen Bitwert}.
  Die Kollapsbedingung ist rein geometrisch und in jeder
  Bit-Interpretation dieselbe. Numerische Gegenprobe:
  $\nthr \cdot \mbit = \Mcoll$ auf neun Stellen (Skript~[2]).
\end{tcolorbox}

\subsection{Die Schwelle}

\begin{equation}
  \nthr = \frac{\Mcoll}{\mbit}
        = \frac{\mP/\sqrt{2}}{\mP\ln 2/(2\pi)}
        = \frac{2\pi}{\sqrt{2}\,\ln 2}
        = 6{,}4097 .
\end{equation}

Die Zahl ist damit ein \emph{Quotient} aus einem geometrischen Zähler
und einem gesetzten Nenner. Alles, was an ihr nicht geometrisch ist,
steckt im Bitwert.

% ============================================================
\section{Dieselbe Rechnung mit systemabhängiger Bit-Energie}

FFGFT setzt statt eines universellen Bitwerts (Dok.~257, A271--A273):
\begin{equation}
  \Ebit(L) = \frac{\hbar c}{L}
  \quad\Longrightarrow\quad
  \mbit(L) = \frac{\hbar}{cL}.
\end{equation}
Eingesetzt in $\nthr = \Mcoll/\mbit$:
\begin{equation}
  \nthr(L) = \frac{\mP}{\sqrt{2}} \cdot \frac{cL}{\hbar}
           = \frac{L}{\sqrt{2}\,\lP},
\end{equation}
da $\lP = \hbar/(c\,\mP)$. Numerisch bestätigt für
$L \in \{1, 5, 10\}\,\lP$ auf Maschinengenauigkeit (Skript~[3]). \B

\begin{tcolorbox}[warnbox, title={Was daraus folgt}]
  $\nthr$ ist eine \textbf{lineare Funktion der Skala}, keine Konstante.
  Sie hat keinen ausgezeichneten Wert, solange keine Skala ausgezeichnet
  ist. Die Frage \enquote{wie viele Bit bis zum Kollaps?} ist ohne
  Angabe der Skala nicht beantwortbar -- sie ist unterbestimmt.
\end{tcolorbox}

\subsection{Matzkes implizite Skala}

Löst man $\mbit = \hbar/(cL)$ nach $L$ auf und setzt Matzkes Bitwert ein:
\begin{equation}
  L_{\mathrm{Matzke}} = \frac{\hbar}{c\,\mbit}
                      = \frac{2\pi\,\lP}{\ln 2}
                      = 9{,}0647\,\lP
                      = 1{,}4651\times10^{-34}~\mathrm{m}.
\end{equation}
Und in der Tat gibt Gl.~\eqref{eq:main} bei dieser Skala
$\nthr = 6{,}4097$ zurück -- exakt Matzkes Wert (Skript~[4]). \B

\begin{tcolorbox}[ffgftbox, title={Kernbefund}]
  Matzkes universeller Bitwert ist die FFGFT-Bit-Energie
  $\Ebit = \hbar c/L$, ausgewertet bei $L = 2\pi\lP/\ln 2$.
  Beide Frameworks rechnen mit derselben Größe; der Unterschied ist,
  dass FFGFT die Skala als freien Parameter mitführt, während Matzke
  eine bestimmte Skala fixiert und das Ergebnis universell nennt.
\end{tcolorbox}

% ============================================================
\section{Numerische Spannweite}

$\nthr$ an den Skalen des Korpus (Skript~[5]):

\begin{center}
\small
\begin{tabular}{@{}llr@{}}
  \toprule
  \textbf{Skala} & $L$ [m] & $\nthr$ \\
  \midrule
  $L_0 = \xi\lP$ (FFGFT-Minimalskala, Dok.~180) & $2{,}155\times10^{-39}$ & $9{,}43\times10^{-5}$ \\
  $\lP$ (Plancklänge)                            & $1{,}616\times10^{-35}$ & $0{,}7071$ \\
  $2\pi\lP/\ln 2$ (Matzkes implizite Skala)      & $1{,}465\times10^{-34}$ & $\mathbf{6{,}4097}$ \\
  $1$~fm (Kernphysik)                            & $1{,}000\times10^{-15}$ & $4{,}38\times10^{19}$ \\
  $L_e = \hbar c/E_0$ (Dok.~257)                 & $2{,}685\times10^{-14}$ & $1{,}17\times10^{21}$ \\
  $1$~nm (Atomphysik)                            & $1{,}000\times10^{-9}$  & $4{,}38\times10^{25}$ \\
  \bottomrule
\end{tabular}
\end{center}

Spannweite über diese Skalen: Faktor $4{,}64\times10^{29}$.
Eine Größe, die über 29 Größenordnungen variiert, ist keine Konstante. \B

\subsection{Empfindlichkeit gegen die Temperaturwahl}

Gleichwertig lässt sich die Abhängigkeit über die angesetzte Temperatur
ausdrücken (Skript~[6]); $\nthr$ skaliert exakt invers zu ihr:

\begin{center}
\small
\begin{tabular}{@{}lrr@{}}
  \toprule
  Angesetzte Temperatur & $\mbit/\mP$ & $\nthr$ \\
  \midrule
  $0{,}5\,T_{\mathrm{KMS}}$ & $0{,}0552$ & $12{,}82$ \\
  $0{,}9\,T_{\mathrm{KMS}}$ & $0{,}0993$ & $7{,}12$ \\
  $1{,}0\,T_{\mathrm{KMS}}$ & $0{,}1103$ & $\mathbf{6{,}41}$ \\
  $1{,}1\,T_{\mathrm{KMS}}$ & $0{,}1213$ & $5{,}83$ \\
  $2{,}0\,T_{\mathrm{KMS}}$ & $0{,}2206$ & $3{,}20$ \\
  \bottomrule
\end{tabular}
\end{center}

Bereits eine Abweichung von $10\,\%$ in der angesetzten Temperatur
verschiebt $\nthr$ um $10\,\%$ -- und damit $\mathrm{Cl}(6)$ von
$0{,}41$~Bit unterhalb der Schwelle auf $1{,}12$~Bit oberhalb
(bei $1{,}1\,T_{\mathrm{KMS}}$: $\nthr = 5{,}83 < 6{,}000$).
Die Stabilitätsaussage kippt.

% ============================================================
\section{Einordnung}

\subsection{Was fällt}

Die Deutung von $\nthr = 6{,}41$ als \emph{fundamentale Zahl} des
Universums. Sie ist der Quotient aus einer geometrischen Größe und einem
gesetzten Bitwert; die Setzung ist eine Skalenwahl, und über die Skalen
variiert das Ergebnis um 29 Größenordnungen. Die Aussage
\enquote{$\mathrm{Cl}(6)$ liegt $0{,}41$~Bit unter der Schwelle}
gilt nur bei $L = 9{,}06\,\lP$ und hat keinen Bestand darüber hinaus. \X

\subsection{Was bleibt}

Erheblich mehr, als die Kritik vermuten lässt:

\begin{itemize}
  \item $\Mcoll = \mP/\sqrt{2}$ aus Compton~$=$~Schwarzschild -- rein
    geometrisch, unabhängig von jeder Bit-Interpretation \B.
  \item Der Schwarzschild-Radius ohne Feldgleichungen: Matzkes
    Novel Result \#1 ist von dieser Analyse unberührt \B.
  \item $\mathrm{Cl}(6)$ als Fermion-Generation (Furey), das
    Herausfallen der Drei, die algebraische Struktur insgesamt \B.
  \item Der Hawking-Mechanismus als algebraisches Bild -- die
    strukturellen Korrespondenzen zu FFGFT sind exakt (Dok.~326).
\end{itemize}

\subsection{Konsequenz für FFGFT}

Die Analyse ist keine bloße Widerlegung; sie stützt die
FFGFT-Position aus Dok.~257 und A271--A273 von außen.
Wenn ein Framework, das explizit einen universellen Bitwert ansetzt,
sich als Auswertung der skalenabhängigen Formel $\Ebit = \hbar c/L$
bei einer bestimmten Skala rekonstruieren lässt, dann ist die
Skalenabhängigkeit nicht der Sonderfall, sondern die allgemeine Form.

Das fügt sich zu Dok.~325: Dort ist die Hawking-Temperatur ohne
universellen Bitwert herleitbar
($k_B T_H = \Ebit(4\pi r_s)$ exakt), und die numerische Gegenprobe
in Dok.~326 ergab $T_H(\text{Matzke})/T_H(\text{KMS}) = 1{,}000000$.
Beide Befunde sagen dasselbe: Der universelle Bitwert ist eine
Umparametrisierung, keine unabhängige physikalische Eingabe.

Bemerkenswert ist dabei die Herkunft des $2\pi$ in Matzkes Bitwert:
Es stammt aus der KMS-Periodizität, also aus genau der Relation, die
in Dok.~325 die Hawking-Temperatur \emph{ohne} Bitwert erzeugt.

% ============================================================
\section{Statusbilanz}

\begin{center}
\begin{tabular}{@{}lc@{}}
  \toprule
  \textbf{Aussage} & \textbf{Status} \\
  \midrule
  $\mbit = \mP\ln 2/(2\pi)$ aus Landauer bei $\TP/(2\pi)$ reproduziert & \B \\
  $\Mcoll = \mP/\sqrt{2}$ rein geometrisch, bitwertfrei & \B \\
  $\nthr(L) = L/(\sqrt{2}\lP)$ -- linear in der Skala & \B \\
  Matzkes Skala: $L = 2\pi\lP/\ln 2 = 9{,}0647\,\lP$ & \B \\
  Spannweite über Korpus-Skalen: Faktor $4{,}6\times10^{29}$ & \B \\
  $\nthr \propto 1/T$ -- $10\,\%$ Temperaturfehler kippt die Aussage & \B \\
  \midrule
  $\nthr = 6{,}41$ als universelle Zahl & \X \\
  Compton~$=$~Schwarzschild, $\mathrm{Cl}(6)$-Struktur, $r_s$ ohne ART & \B\ (unberührt) \\
  \bottomrule
\end{tabular}
\end{center}

% ============================================================
\section{Vorgeschlagener Registereintrag R88}

\textbf{R88 -- $\nthr$ ist skalenabhängig; Matzkes Skala rekonstruiert
(August 2026):}
Dok.~329 entscheidet die in Dok.~326 als \S\ geführte Frage numerisch.
Die Kollapsbedingung Compton~$=$~Schwarzschild liefert
$\Mcoll = \mP/\sqrt{2}$ rein geometrisch und bitwertfrei \B.
Mit der systemabhängigen Bit-Energie $\Ebit(L) = \hbar c/L$ (Dok.~257)
folgt exakt $\nthr(L) = L/(\sqrt{2}\lP)$ -- linear in der Skala \B.
Matzkes Wert $6{,}4097$ entspricht exakt
$L = 2\pi\lP/\ln 2 = 9{,}0647\,\lP$ \B; der \enquote{universelle}
Bitwert ist die FFGFT-Bit-Energie bei dieser Skala.
Über die Skalen des Korpus variiert $\nthr$ um Faktor $4{,}6\times10^{29}$;
eine Temperaturabweichung von $10\,\%$ kippt die Stabilitätsaussage für
$\mathrm{Cl}(6)$. Damit ist $\nthr = 6{,}41$ als universelle Zahl
abgeschlossen negativ \X.
Unberührt bleiben Compton~$=$~Schwarzschild, der Schwarzschild-Radius
ohne Feldgleichungen und die $\mathrm{Cl}(6)$-Struktur \B.
Prüfskript: \texttt{326\_nthresh\_skalenabhaengigkeit.py} (12 Assertionen).\\
\emph{Deklariert 24.~August 2026}

% ============================================================
\section*{Bezug zum Korpus}

\begin{center}
\begin{tabular}{@{}ll@{}}
  \toprule
  \textbf{Baustein} & \textbf{Dok.} \\
  \midrule
  Bit-Energie $\Ebit = \hbar c/L$, Informationseinheit $\Delta w = 1$ & 257 \\
  Landauer als Träger-Operation, nicht abstrakte Information & A271--A273 \\
  Elementarzellen, Bitkapazität, Partitionsinvarianz & 302 \\
  $\lP$, $\mP$, $L_0 = \xi\lP$ & 180 \\
  KMS-Relation, $k_BT_H = \Ebit(4\pi r_s)$ & 325 \\
  Matzke-Vergleich, Bitwert-Divergenz & 326 \\
  \bottomrule
\end{tabular}
\end{center}

\section*{Literatur}

\begin{enumerate}
  \item D.~Matzke, \textit{Black Holes from Hyperbits},
    IPI Annual Conference, August 2026.
  \item J.~Pascher, Dok.~326 \textit{Schwarze Löcher aus Hyperbits und
    aus FFGFT}, FFGFT-Korpus, 2026.
  \item J.~Pascher, Dok.~325 \textit{Der FFGFT-Hawking-Mechanismus},
    FFGFT-Korpus, 2026.
  \item J.~Pascher, Dok.~257, 302, A271--A273 (Bit-Energie, Landauer),
    FFGFT-Korpus.
  \item J.~Pascher, Dok.~180 \textit{Herleitung von $L_0$},
    FFGFT-Korpus, 2026.
  \item R.~Landauer, IBM J.\ Res.\ Dev.\ \textbf{5}, 183, 1961.
\end{enumerate}

\vspace{6pt}\noindent\small
\textcopyright\ 2026 Johann Pascher ·
\href{https://creativecommons.org/licenses/by/4.0/}{CC BY 4.0}
